\chapter*{Abstract}
La presente tesi ha l’obiettivo di fornire una panoramica sull’utilizzo degli algoritmi genetici (GA) per l’addestramento delle reti neurali. In particolare, il lavoro si concentra sul caso delle Multi-Layer Perceptron (MLP), la cui architettura verrà approfondita nei capitoli successivi.

Verranno analizzati i principali approcci che impiegano i GA in questo contesto, a partire dalla codifica di una rete come individuo e dai meccanismi di generazione della nuova popolazione ispirati alla teoria darwiniana.

Nella parte più tecnica saranno presentati alcuni esperimenti, con l’obiettivo di evidenziare le problematiche tipiche dei GA, in particolare l’elevato costo computazionale e i tempi richiesti per raggiungere prestazioni ottimali. Saranno inoltre discusse alcune strategie volte a ridurre, almeno parzialmente, tali criticità.

Infine, verrà analizzato l’aspetto forse più interessante: l’impiego dei GA in combinazione con i metodi classici basati sul gradiente. In particolare, si evidenzierà come l’approccio ibrido possa contribuire a superare il problema dei minimi locali, caratteristico delle tecniche di ottimizzazione gradient-based.